% 391_Puzzleteile_Kosmologie_En_ch.tex
% Johann Pascher, 9 Oct 2026
% Loose Puzzle Pieces II: Cosmology

\providecommand{\BM}{\textbf{[B]}}
\providecommand{\KM}{\textbf{[K]}}
\providecommand{\SM}{\textbf{[S]}}
\providecommand{\XM}{\textbf{[X]}}
\providecommand{\QM}{\textbf{[Q]}}
\providecommand{\EM}{\textbf{[E]}}
\providecommand{\SetzM}{\textbf{[SETZUNG]}}

\chapter*{Loose Puzzle Pieces II: Cosmology}
\addcontentsline{toc}{chapter}{Loose Puzzle Pieces II: Cosmology}

\begin{abstract}
Doc.~390 placed the building blocks of particle physics that were found in the literature afterwards
and fit into the picture of FFGFT. This document does the same for cosmology. It goes through eight
building blocks: Zwicky's tired light, the gravitation theory of Hoyle and Narlikar and Wetterich's
universe without expansion, Einstein's static universe and the steady-state theory, Dirac's large
numbers, the Hubble length, Milgrom's acceleration scale, the problem of vacuum energy, and the
cosmic background radiation. The order holds here too: the building blocks were not the foundation
of FFGFT; they fit in afterwards. Unlike in particle physics, however, the cosmic sector lies outside
the claim of derivation: $H_0$ is calibrated to the measurement, the exponent~10 is a setting, and
the Standard Model does not derive this sector either. The puzzle pieces therefore show readings,
not derivations: what the observations of cosmology carry can be read with $\tilde T\cdot m=1$
without expansion, and several places in the literature have already shown this. Two building
blocks, vacuum energy and background radiation, do not fit so far or fit only as a candidate, and
are named as the boundary.
\end{abstract}

% ============================================================
\section*{Status markers}
% ============================================================
\begin{center}
\begin{tabular}{@{}lp{11cm}@{}}
\textbf{[SETZUNG]} & Motivated input, not a derivation. \\[2pt]
\textbf{[B]} & Algebraically proved (verification script, exact arithmetic). \\[2pt]
\textbf{[K]} & Numerically confirmed. \\[2pt]
\textbf{[S]} & Structurally plausible, identification open. \\[2pt]
\textbf{[E]} & Established mathematics or physics, adopted. \\[2pt]
\textbf{[Q]} & Taken from the literature, with source. \\[2pt]
\textbf{[X]} & Checked and negative. \\
\end{tabular}
\end{center}
The markers give the status \emph{in the corpus}; the document derives nothing new but places things
in context. Verification script: \texttt{pruef\_391\_kosmologie.py}.

% ============================================================
\section{Order and boundary}
% ============================================================

\textbf{On the order.} As in Doc.~390, none of the following building blocks was a starting point of
FFGFT. The theory was developed from $\tilde T\cdot m=1$ and the geometry of the carrier; the
literature references were found only afterwards. Where it says that a building block fits in at
some place in the corpus, this is agreement in hindsight, not origin.

\textbf{On the boundary.} The cosmic sector is deliberately left out of the corpus. The Standard
Model does not derive $H_0$ and $\Lambda$ either, so a comparison would be asymmetric (P39 in
Doc.~190; Doc.~365). The form of the relation between the Hubble length and the Planck length is
fixed as a dimensionless $\xi$ ratio (P20), the exponent~10 in $H_0\propto\xi^{10}$ is a setting
\SetzM, and the scale is calibrated to the measurement of $H_0$ (R137 in Doc.~190). What this
document places are therefore readings of the observations, not predictions of their values.

% ============================================================
\section{Redshift without expansion}
% ============================================================

\subsection{Zwicky's tired light (1929)}

\textbf{What it contained.} Shortly after Hubble's discovery, Zwicky proposed that the redshift of
distant galaxies arises not from expansion but from light losing energy along the way
\QM~\cite{Zwicky1929}. This was the first static interpretation of the Hubble relation.

\textbf{What was missing.} Tired light changes only the photon. Distant clocks, such as the light
curves of supernovae, would then have to run at their original rate. What is observed, however, is
a stretching by $(1+z)$, and the idea failed on that.

\textbf{Where it fits in the corpus.} Because of $\tilde T\cdot m=1$, wavelength and clock rate are
booked together. From $1+z=e^{\xi_0 x}$ the time dilation $(1+z)$ follows for all $z$ \BM{} (R128 in
Doc.~190). Classical tired light fails because it changes only the photon side; the mass reading
cannot make this mistake structurally (Doc.~312). The redshift is achromatic, and the scale of the
path lengthening is calibrated via $H_0$ (K6 in Doc.~190).

\subsection{Hoyle--Narlikar (1964) and Wetterich's universe without expansion (2013)}

\textbf{What it contained.} Hoyle and Narlikar formulated a theory of gravitation in which particle
masses arise from the distribution of matter in the cosmos \QM~\cite{HoyleNarlikar1964}. In 2013
Wetterich showed that an expanding universe with constant masses turns, through a field
redefinition, exactly into a static universe with growing particle masses \QM~\cite{Wetterich2013}.
Only dimensionless ratios are observable.

\textbf{What was missing.} A reason why the running of masses should be the natural reading and not
merely an equivalent rewriting.

\textbf{Where it fits in the corpus.} Doc.~312 lists the equivalence of expansion and running masses
as \KM{} and thereby anchors the degeneracy from Doc.~267: redshift, luminosity distance, baryon
acoustic oscillations and background temperature do not distinguish the two readings. In FFGFT the
running of masses is not an additional assumption, because mass and clock rate are the same quantity
through $\tilde T\cdot m=1$ anyway. Doc.~312 records the boundary: the equivalence shows that the
static reading is possible and observationally equal, not that it is right.

\subsection{Einstein's static universe (1917) and the steady-state theory (1948)}

\textbf{What it contained.} In 1917 Einstein built the first relativistic world model as a static,
spatially closed cosmos \QM~\cite{Einstein1917}. In 1948 Bondi, Gold and Hoyle formulated a universe
without a beginning that looks the same on large scales at all times
\QM~\cite{BondiGold1948,Hoyle1948}.

\textbf{What was missing.} Einstein's model needed a specially introduced cosmological constant and
was unstable. The steady-state theory needed a continuous creation of matter, and the discovery of
the background radiation in 1965 was taken as its end.

\textbf{Where it fits in the corpus.} FFGFT lives on a carrier without boundary,
$\partial T^4=\emptyset$ (Doc.~312). Doc.~313 works through the question of the beginning on the
compact time cycle: the observable history fills half a cycle, and the \enquote{beginning} is the
antipode of the cycle, not an event \KM{} with the calibration from P20. Neither a cosmological
constant nor creation of matter is needed for this. According to Doc.~313 it remains open that an
exact periodicity of the cycle is not dynamically generic and has to be set as a boundary condition.

% ============================================================
\section{Scales}
% ============================================================

\subsection{Dirac's large numbers (1937)}

\textbf{What it contained.} Dirac noticed that the ratio of the electric to the gravitational force
between electron and proton and the age of the universe in atomic units of time are both about
$10^{39}$, and suspected a connection behind this \QM~\cite{Dirac1937}.

\textbf{What was missing.} A mechanism. Dirac's conclusion that the gravitational constant must
decrease with time is excluded by measurements.

\textbf{Where it fits in the corpus.} In the corpus the relation between Hubble length and Planck
length is a dimensionless $\xi$ ratio (P20); the exponent is set, and the scale is calibrated to
$H_0$. P39 records that no framework derives this relation from first principles. Dirac's
observation thus appears as a ratio of two scales of the same geometry, not as a constant running in
time. Dirac's name does not yet occur in the corpus; the placement is a reading \SM.

\subsection{The Hubble length as maximal scale}

\textbf{What it contained.} In 1929 Hubble found the linear relation between distance and redshift
\QM~\cite{Hubble1929}. Since then the length $c/H_0$ has been read as the horizon of an expanding
universe.

\textbf{What was missing.} A reading of this length without expansion.

\textbf{Where it fits in the corpus.} Doc.~182 and Doc.~365 read $R_H=c/H_0$ as the geometric maximal
scale of the compact carrier, not as an expansion horizon. Its value is calibrated to $H_0$.

\subsection{Milgrom's acceleration scale (1983)}

\textbf{What it contained.} Milgrom explained the flat rotation curves of galaxies without dark
matter through an acceleration scale $a_0$ below which Newton's law deviates \QM~\cite{Milgrom1983}.
He already noticed that $a_0$ is of the order of $cH_0$.

\textbf{What was missing.} A reason for this closeness to the cosmic scale, and a justification of
the interpolation formula.

\textbf{Where it fits in the corpus.} In the corpus $a_0=cH_0/(2\pi)$ \KM{} (Doc.~365); the same
exponent~10 that carries $H_0$ also reproduces $a_0$, which would not be expected from a pure $H_0$
fit (Doc.~309, cross-anchoring). Doc.~308 obtains the form $a=\sqrt{a_N^2+a_N a_0}$ with
$a_0=2a_{\mathrm{bg}}$ from the Unruh coupling \KM/\SM. The book \enquote{Dark Matter -- an
Illusion} works out the rotation curves. The exponent remains set.

% ============================================================
\section{Where it does not fit yet, or only as a candidate}
% ============================================================

\subsection{Vacuum energy and the cosmological constant}

\textbf{What it contained.} The zero-point energy of quantum fields is measured in the Casimir
effect. As a source of the cosmological constant, its naive estimate lies many orders of magnitude
above the observed value; Weinberg called this the cosmological constant problem
\QM~\cite{Weinberg1989}.

\textbf{Status in the corpus.} The connection between the Casimir and background energy densities at
the length $L_\xi$ is an identity, not evidence from measurement \SM{} (R136 in Doc.~190). The
$\Lambda$ formulas of older documents do not hit the observed value \XM{} (R137). In the static
reading, $\Lambda$ is an artefact of the reading without a referent of its own (Doc.~308) and is not
set (Doc.~312); there is no derivation of the observed value. This building block does not fit in
so far.

\subsection{The cosmic background radiation (1965)}

\textbf{What it contained.} In 1965 Penzias and Wilson found an isotropic radiation of about
$3$~kelvin \QM~\cite{PenziasWilson1965}. Its spectrum and its anisotropies are among the most precise
data in cosmology.

\textbf{Status in the corpus.} The relations for $T_{\mathrm{CMB}}$ from $\xi$ give a pure number
that agrees with the measured value only in the unit eV; they count as a unit-dependent agreement,
and the question is open (R137, R70). The relation $T(z)$ in Doc.~039 is excluded by measurements
\XM. The factor~3 in front of the peak sequence is a setting, and the peak selection as a whole is
open (R139). In the basic form, Doc.~388 sets up the candidate
$T_{\mathrm{CMB}}/m_e=(8\pi)^{1/4}\,\xi^{5/2}$, which hits the FIRAS value \KM; the prefactor is
motivated but not justified \SM, and the random test there shows that chance is not excluded. If the
candidate holds, it connects temperature and Hubble scale through $T_{\mathrm{CMB}}^4=16\,H_0\,m_e^3$
\BM{} (Doc.~388), with the same set exponent~10. The temperature of the radiation is moreover among
the quantities that do not resolve the degeneracy from Doc.~267. This building block thus fits in
only as a candidate, not as a secured place.

% ============================================================
\section{What the pieces show}
% ============================================================

In particle physics (Doc.~390) the building blocks fitted in at places where the corpus derives
something. In cosmology they fit in at places where the corpus \emph{reads} something. Zwicky, Hoyle
and Narlikar, Wetterich, Einstein and the steady-state theory each supplied a piece of a static
cosmology; each piece failed on its own at a different point, at time dilation, at the missing
justification of running masses, at instability or at creation of matter. With $\tilde T\cdot m=1$
these points are no longer separate problems, because wavelength, clock rate and mass are the same
booking.

Dirac and Milgrom each noticed a striking closeness between microscopic and cosmic scales. In the
corpus both are ratios of the same geometry, but tied to a set exponent and a calibrated scale. They
are therefore cross-anchorings, not evidence.

Vacuum energy and background radiation do not fit so far, or only as a candidate. This is not
concealed: the formulas concerned are listed in the register as \XM{} or as open, and the temperature
candidate from Doc.~388 as \SM. For the cosmic sector the same therefore holds as for
scheme-dependent quantities in particle physics: no claim of derivation is made where the quantities
depend on the measurement or on a convention.

% ============================================================
\section{Overview}
% ============================================================

\begin{center}
\small
\begin{tabular}{@{}>{\raggedright\arraybackslash}p{3.0cm}>{\raggedright\arraybackslash}p{1.7cm}>{\raggedright\arraybackslash}p{4.1cm}>{\raggedright\arraybackslash}p{3.8cm}@{}}
\toprule
Building block & Year & What was missing & In the corpus \\
\midrule
Zwicky, tired light & 1929 & no time dilation & R128, K6; Doc.~312 \\
Hoyle--Narlikar, Wetterich & 1964, 2013 & running masses only equivalent, not justified & Doc.~312, 267 \\
Einstein universe, steady state & 1917, 1948 & instability, creation of matter & Doc.~313, 312 \\
Dirac's large numbers & 1937 & no mechanism & P20, P39 \SM \\
Hubble length & 1929 & read only as expansion horizon & Doc.~182, 365 \\
Milgrom's $a_0$ & 1983 & closeness to $cH_0$ unexplained & Doc.~365, 309, 308 \\
Vacuum energy, $\Lambda$ & 1989 & does not fit so far & R136, R137 \\
Background radiation & 1965 & only a candidate, chance not excluded & Doc.~388; R137, R139, R70 \\
\bottomrule
\end{tabular}
\end{center}

% ============================================================
\section{Conclusion}
% ============================================================

In cosmology, too, the puzzle pieces were in the literature long before FFGFT, and here too they were
not its foundation. Six of the eight building blocks fit in afterwards, as readings of the
observations without expansion and as cross-anchorings between microscopic and cosmic scales. Two do
not fit so far or fit only as a candidate. The difference from particle physics must be stated
honestly: the cosmic sector is calibrated and left out; here the building blocks show that the
theory can read cosmology, not that it derives its numbers.

\vspace{1em}
\noindent\textit{Verification script:} \texttt{2/python/Dok391\_Skripte/pruef\_391\_kosmologie.py}
--- 51/51.

% ============================================================
\section*{Tool and method}
\addcontentsline{toc}{section}{Tool and method}
% ============================================================

Claude (Anthropic) was used as a tool for: compiling the references in the corpus, writing and
running the verification script, and working out the \LaTeX{} sources according to the author's
specifications. The verification script backs every corpus reference in this document with a text
check in the named chapter file; if a check fails, the reference is not used.

\begin{thebibliography}{99}

\bibitem{Zwicky1929}
F.~Zwicky, \textit{On the red shift of spectral lines through interstellar space},
Proc.\ Natl.\ Acad.\ Sci.\ USA \textbf{15} (1929) 773--779.

\bibitem{HoyleNarlikar1964}
F.~Hoyle, J.~V.~Narlikar, \textit{A new theory of gravitation},
Proc.\ R.\ Soc.\ Lond.\ A \textbf{282} (1964) 191--207.

\bibitem{Wetterich2013}
C.~Wetterich, \textit{A Universe without expansion},
Phys.\ Dark Univ.\ \textbf{2} (2013) 184--187; arXiv:1303.6878.

\bibitem{Einstein1917}
A.~Einstein, \textit{Kosmologische Betrachtungen zur allgemeinen Relativitätstheorie},
Sitzungsber.\ Preuß.\ Akad.\ Wiss.\ Berlin (1917) 142--152.

\bibitem{BondiGold1948}
H.~Bondi, T.~Gold, \textit{The steady-state theory of the expanding universe},
Mon.\ Not.\ R.\ Astron.\ Soc.\ \textbf{108} (1948) 252--270.

\bibitem{Hoyle1948}
F.~Hoyle, \textit{A new model for the expanding universe},
Mon.\ Not.\ R.\ Astron.\ Soc.\ \textbf{108} (1948) 372--382.

\bibitem{Dirac1937}
P.~A.~M.~Dirac, \textit{The cosmological constants}, Nature \textbf{139} (1937) 323.

\bibitem{Hubble1929}
E.~Hubble, \textit{A relation between distance and radial velocity among extra-galactic nebulae},
Proc.\ Natl.\ Acad.\ Sci.\ USA \textbf{15} (1929) 168--173.

\bibitem{Milgrom1983}
M.~Milgrom, \textit{A modification of the Newtonian dynamics as a possible alternative to the
hidden mass hypothesis}, Astrophys.\ J.\ \textbf{270} (1983) 365--370.

\bibitem{Weinberg1989}
S.~Weinberg, \textit{The cosmological constant problem}, Rev.\ Mod.\ Phys.\ \textbf{61} (1989) 1--23.

\bibitem{PenziasWilson1965}
A.~A.~Penzias, R.~W.~Wilson, \textit{A measurement of excess antenna temperature at 4080 Mc/s},
Astrophys.\ J.\ \textbf{142} (1965) 419--421.

\end{thebibliography}
